\section{Discussion}\label{sec:discussion}

\subsection{What the benchmark teaches about multi-source yield forecasting}
Three lessons stand out. The first is that \emph{leakage-safe evaluation must reach every loop}. The most consequential failure in this study did not arise from the outer protocol, which was leakage-safe, but from an inner routine---the efficient leave-one-out selection of the Ridge penalty---that silently treated district-seasons as exchangeable. Within-season error correlation let season-level predictors such as prices and reservoir storage fit season effects, the penalty was set too small, and the model extrapolated badly when a new season brought new values of these predictors. Random outer cross-validation did not reveal the problem; it reported excellent and stable performance precisely because it never asked the model to forecast an unseen season. The practical rule is simple: every resampling step, inner and outer, should be blocked by season (and by space when forecasting unseen districts) \citep{roberts2017cross,bergmeir2012use}.

The second lesson is that \emph{stability is a first-class model property}. Tree ensembles extrapolate conservatively, holding predictions constant beyond the training range, whereas linear models extrapolate along every predictor. With fewer than fifteen training seasons and several predictors that vary mainly between seasons, this asymmetry dominated: the tuned LightGBM varied by only \repSdLGB~\kgrai{} across ten worlds, while Ridge failed in \repFailRidge{} of them. Ensembles inherit the failures of their members unless each member is individually robust; once Ridge was tuned by season, the ensemble became the most accurate learner (an exploratory result). The neural baseline did not beat the tree models, in line with benchmarks on tabular data \citep{grinsztajn2022why,shwartz2022tabular}.

The third lesson is that \emph{explanations need ground truth}. Grouped SHAP values reflected strong, well-observed mechanisms faithfully but misattributed correlated mechanisms (heat to moisture) and season-level effects (trend to prices), consistent with known difficulties of attributions under dependence \citep{aas2021explaining,hooker2021unrestricted}. Explanations also inherit model fragility: the ensemble attributed the largest share of its forecasts to prices, although removing prices improved it and its price attributions ran against the true response. In an empirical study such attributions could easily be read as evidence of economic effects. Grouped attributions of stable models, cross-checked against ablation and, where possible, against simulations with known truth, are safer; causal questions about prices require identification strategies from econometrics rather than predictive attributions \citep{wooldridge2010econometric}.

\subsection{Implications for operational forecasting}
If the benchmark's structure is representative of the region, a district-level system built from MODIS-class composites, ERA5-Land and CHIRPS weather, reservoir data and official statistics could issue informative forecasts from July and sharpen them substantially in October, when satellite information starts to add value, with 80\% intervals about \widthEighty~\kgrai{} wide. Such a system would serve the agencies that compile yield statistics, monitor floods from space, manage reservoirs and administer crop insurance. The leakage audit implies that accuracy figures obtained with random cross-validation should not be used to justify such systems, and the failure analysis implies that model stability across seasons deserves as much attention as average accuracy.

\subsection{Limitations}
All conclusions are conditional on the DGP. Although its mechanisms are drawn from the agronomic and remote-sensing literature, the DGP is additive with hinge-type stresses, its parameters were set to plausible values rather than estimated from Thai data, and its satellite indices are generated from the visible part of yield; real indices also respond to weeds, mixed pixels, varietal and phenological differences, so the benchmark may overstate the value of satellite data. Only eight test seasons were available, which limits the power of season-level tests; replicate worlds address Monte-Carlo variability but not structural misspecification. The configuration was tuned once in the main world and reused elsewhere. SHAP values were computed in their interventional form; conditional variants may behave differently under dependence. The analysis is at district level and does not use field-level or radar data, and spatial cross-validation for districts without yield history was not evaluated. Finally, an exploratory version of this pipeline was run before the final design: although tuning and model selection used validation seasons only, the analyst was not blind to earlier test-period results, and the diagnosis of the linear failures is post hoc. Moreover, the post-hoc check was applied to worlds selected because they had failed; a prospective test of season-blocked tuning on fresh worlds is required before the remedy can be recommended on empirical grounds. Recovery was measured by correlation, which ignores the scale of attributions, and the neural baseline was not tuned.

\subsection{Pathway to real data}
The code ingests real inputs with the same schema (Table~\ref{tab:features}): yield statistics, CHIRPS and ERA5-Land weather, MODIS composites extracted with the provided Earth Engine script, farm-gate prices, reservoir storage and national soil, land-use and flood maps. Natural extensions are Sentinel-1 radar to observe the crop through monsoon clouds, Sentinel-2 for field-level models, spatially blocked validation for districts without history, conformalised quantile regression or adaptive conformal inference for heteroscedastic and shifting errors \citep{romano2019conformalized,gibbs2021adaptive}, and season-blocked inner tuning for every learner.

\section{Conclusion}\label{sec:conclusion}
In a benchmark with known truth, multi-source machine learning reduced district-level rice yield forecast errors by \redClimPct--\redClimLGBPct\% relative to climatology, with satellite data as the indispensable predictor group and with most of their value arriving from October. The same benchmark exposed three pitfalls that real-data studies cannot see: random cross-validation overstated skill and hid failures; leakage in an inner tuning loop caused catastrophic failures of a linear learner; and SHAP explanations were faithful only for strong, well-observed mechanisms and stable learners. We therefore recommend (i) blocking every resampling step by season, (ii) reporting skill relative to climatology and on anomalies, (iii) preferring learners that are stable under season-level covariate shift, or tuning linear components by season---a remedy our exploratory analysis supports but that still needs prospective confirmation, (iv) treating attributions as descriptions of a model and validating them by simulation, and (v) reporting calibrated intervals together with their coverage by season. The benchmark and code are released so that these checks can precede the use of real data.
